\section{Introduction}\label{sec:introduction}
All graphs in this paper are finite, simple, and undirected. A
\emph{cycle--edge decomposition} is a partition of the original edge set
into edge sets of simple cycles and singleton edges. Let $\ce(G)$ be the
minimum number of parts in such a partition. We write
\[
 n=\abs{V(G)},\qquad m=\abs{E(G)},\qquad r(G)=n-c(G),
\]
where $c(G)$ counts all connected components, including isolated vertices.
The empty graph on no vertices has $c(G)=0$.

\begin{theorem}\label{thm:main}
Every finite simple graph $G$ satisfies
\[
                 \ce(G)\le 12\bigl(\abs{V(G)}-c(G)\bigr).
\]
\end{theorem}

The problem of obtaining a linear bound for $\ce(G)$ goes back to Erd\H{o}s
and Gallai; it appears in the discussion of Erd\H{o}s, Goodman and
P\'osa~\cite{EGP1966}. Conlon, Fox and Sudakov~\cite{CFS2014} proved an
$O(n\log\log n)$ bound, and Buci\'c and Montgomery~\cite{BM2024} subsequently
obtained $O(n\log^*n)$. Linear bounds were presented in 2026 by
OpenAI~\cite{OAI2026}, by Kim~\cite{Kim2026}, and in a separate manuscript
and formalization of Coffey~\cite{Coffey2026}. The present paper addresses an
explicit coefficient. Its large-scale construction follows the
quotient-and-lifting approach of~\cite{OAI2026}; the result is not a new proof
of the mere existence of some linear constant.

Our proof first assigns weight $1$ to a cycle and weight
$\sigma=23/42$ to a singleton edge. With $C=485/42$, the central estimate is
\begin{equation}\label{eq:weighted-order}
 \exists\mathcal P\quad \wt_\sigma(\mathcal P)\le Cn,
\end{equation}
where $\mathcal P$ is an actual cycle packing and all uncovered original
edges are counted as singletons. A component-rank normalization and the
identity
\[
                  C+1-\sigma=12
\]
then give Theorem~\ref{thm:main}. Keeping the singleton weight throughout
the proof is essential: replacing weighted completion costs by unweighted
ones at intermediate stages loses the final coefficient.

The quantitative proof has three parts. First, a mixed cycle-packing estimate
with constants $43/10$ and $3$ is combined with a short-circumference base
and finite depth-first-search chains. Second, an indexed scale decomposition
uses three final degree-reduction rates and preserves the exact assignment of
every original edge. Third, a positive renewal potential pays the accumulated
losses. Forty stopping-phase intervals determine different prefix thresholds;
a positive second-order recurrence controls every possible number of layers.

The proof is computer-assisted. Two local interfaces require substantial
finite certificates: the short-circumference base and the quantitative
mixed-cover estimate. Their mathematical specifications and certification
procedures are given in Sections~\ref{sec:finite}, \ref{app:terminal}, and
\ref{app:cover}. The subsequent scale and cost calculations are displayed
explicitly, with the rational data and the complete formal proofs supplied
in the companion repository. Appendix~\ref{app:correspondence} records the
correspondence between the principal statements here and their Lean
counterparts. The coefficient throughout this paper is exactly $12$; no
parameters from later numerical experiments are used.
